\documentclass{beamer}

\newcommand{\week}{Financial Markets 1 of 8}

\title{Financial Market Returns}
\subtitle{Mishkin Chapter 4}
\author{Econ 357}
\date{\week}

% Reference the shared preamble
\input{preamble}

% Set the footer -- change 
\setbeamertemplate{footline}{
  \leavevmode%
  \vspace{2ex}
  \hbox{%
    % Left box: Econ 457
    \begin{beamercolorbox}[wd=.4\paperwidth,ht=2.5ex,dp=1ex,left]{author in head/foot}%
      \hspace{1em}Econ 357
    \end{beamercolorbox}%
    % Middle box: Week
    \begin{beamercolorbox}[wd=.2\paperwidth,ht=2.5ex,dp=1ex,center]{date in head/foot}%
      \centering\week
    \end{beamercolorbox}%
    % Right box: Slide numbers
    \begin{beamercolorbox}[wd=.4\paperwidth,ht=2.5ex,dp=1ex,center]{date in head/foot}%
      \hfill\insertframenumber{} 
    \end{beamercolorbox}%
  }%
  \vskip0pt%
}

\begin{document}

\frame{\titlepage}

\begin{frame}[t]
    \frametitle{Review}
        \textbf{Intro 1}
            \begin{itemize}
                \item Real Assets and Financial assets.
                \item Every financial assets is somebody's asset and somebody else's liability.
            \end{itemize}

        \vspace{1em}

        \textbf{Intro 2}
            \begin{itemize}
                \item Financial Intermediaries
                \item Reasons they exist: transaction costs, risk sharing, Asymmetric Information.
                \item Asymmetric Information: Adverse Selection and Moral Hazard.
            \end{itemize}

\end{frame}

\begin{frame}
    \frametitle{Outline for Financial Markets}
        \begin{enumerate}
            \item \fbox{Financial Market Returns}
            \item Bonds 1: Discounting, Prices and Yields
            \item Bonds 2: Nominal v. Real Yields, Supply and Demand
            \item Bonds 3: Market for Money
            \item Bonds 4: Risk and Term Structure
            \item Equities 1: Dividend Discount Model
            \item Equities 2: P/E Ratio
            \item Equities 3: Theories of Stock Pricing
            
        \end{enumerate}
\end{frame}


\begin{frame}
    \frametitle{Outline for Today's Lecture}
    \begin{enumerate}
        \item Financial Market Returns
        \begin{itemize}
            \item Total Return Definition
            \item Measurement
        \end{itemize}
        \item Real v. Nominal Returns
        \item Historical Distribution of Returns
    \end{enumerate}
\end{frame}

\begin{frame}[t]
    \frametitle{1. Financial Market Returns}
    \section{Financial Market Returns}
    \framesubtitle{Total Returns}

    Definition of \textbf{total return}:
    $$\text{Return (\%)} = \frac{\text{Income}}{\text{Initial Price}} + \frac{\text{Change in Price}}{\text{Initial Price}} $$

    \vspace{1em}
    Notes:
    \begin{itemize}
        \item This is a \textit{backwards looking} formula.   
        \item Calculating total return requires three values: beginning and ending price and income.   Given total return and two of the three values, could calculate the third value by rearranging the equation.
        \item Expressed as a percentage of the initial price.
        \item Don't forget income!  Only include income that was earned during the holding period.
    \end{itemize}  
\end{frame}

\begin{frame}[t]
    \frametitle{1. Financial Market Returns}
    \framesubtitle{Income}

    \vspace{-1em}
            \begin{center}
    \textbf{Types of Income}
    \end{center}
    \begin{tblr}{
      colspec = { Q[l,wd=2cm] |Q[c,wd=2cm] |Q[l,wd=6cm,cmd=\scriptsize] }, vline{3} = {1-Z}{}
    }
        \toprule
        Instrument & Income & Notes \\
        \midrule
    Treasury Bond & Coupon &
    Part of bond contract, paid twice per year (semi-annual), amount is \% of bond face value \\
        \midrule

    Corporate Bonds & Coupon &  Part of bond contract, paid twice per year (semi-annual), amount is \% of bond face value \\
        \midrule

    Equities & Dividend & Discretionary (no obligation to pay), paid quarterly, the amount may vary, again at company’s discretion \\
        \midrule

    ETF / Mutual Funds & Distributions & Varies widely by the fund, in many cases reflects coupons and dividends of underlying assets, but could also include capital gains. \\
        \bottomrule
    \end{tblr}
    \vspace{1em}
\end{frame}

\begin{frame}[t]
    \frametitle{2. Income}
    \framesubtitle{Treasury Bond: Note the ``coupons'' at the bottom}
    \centering
    \includegraphics[width=0.5\textwidth]{figures/intro_b_fig3.png}

    \blfootnote{Source: The Joe I. Herbstman Memorial Collection of American Finance}

\end{frame}

\begin{frame}[t]
    \frametitle{2. Income}
    \framesubtitle{Timing of Income}

    Be careful with the timing of income.   
    \begin{itemize}
      \item In order to receive a dividend, you must purchase the stock
    before the ``ex-dividend'' date.
    \item If you purchase a stock on or after the ``ex-dividend'' date,
    you will not receive the upcoming dividend; the seller retains the
    right to the dividend.
    \item A problem for this class may say "you purchase a stock that just paid a dividend." 
    Do NOT count that dividend in the total return, as you purchased it \textit{after} the dividend was paid.
    \end{itemize}

\end{frame}

\begin{practiceframe}[t]
    \frametitle{1. Financial Market Returns}
    \framesubtitle{Measurement}

      Stock X just paid a dividend of \$10.  You then buy Stock X for \$100, hold it for one year during which time it pays 4 more dividends of the same amount.   After one year you then sell the stock for \$110.  What was your total return?

    \centering
    \includegraphics[width=\textwidth]{Econ357/figures/02_mkts_1_timeline.png}

\end{practiceframe}

\begin{frame}[t]
    \frametitle{1. Financial Market Returns}
    \framesubtitle{Measurement}

    Total Returns:
    
    $$
    \text{Return (\%)} = \frac{\text{Income}}{\text{Initial Price}} +
    \frac{\text{Change in Price}}{\text{Initial Price}} $$

    \vspace{2em}
    Alternative terminology:
    
    $$
    \text{Return (\%)} = \text{Income Return} + \text{Capital Gain/Loss}
    $$

\end{frame}

\begin{frame}[t]
    \frametitle{1. Financial Market Returns}
    \framesubtitle{Example: S\&P Total Return}

    \centering
    \includegraphics[width=0.8\textwidth]{Econ457/figures/ch5_1_sp500_tr.png}

    \blfootnote{Data source: Interactive Brokers (SPX and SPXT).   It many cases it is easiest to work with Total Return series, which already include income.   In homework assignment \#1 you will construct a total return series.}

\end{frame}

\begin{practiceframe}[t]
    \frametitle{1. Financial Market Returns}
    \framesubtitle{Measurement}
    
    \begin{enumerate}
    \item Based on the information below, what is the total return for an investor in Coca-Cola (KO) that purchased the stock on Jan 1, 2025 and sold it one year later.   What was the capital gain or loss?   What was the income return?\\
    Price:
        \begin{itemize}
            \item Jan 1, 2025:  \$61.87
            \item Dec 31, 2025: \$69.93
        \end{itemize}
	Dividends: 
        \begin{itemize}
            \item \$0.51 per quarter 
        \end{itemize}

    \item Consider a 10-year Treasury bond that has a 4.25\% coupon, which is paid semi-annually.   An investor purchases this bond for \$97.26 on Jan 1, 2025.   She sells the bond one year later and realizes a total return of 8.2\%.   What price did she sell the bond for?
    \end{enumerate}

\end{practiceframe}


\begin{frame}[t]
    \frametitle{2. Real v. Nominal Returns}
    \section{Real v Nominal Returns}
    \framesubtitle{What Is Inflation?}

    \begin{itemize}
        \item Inflation is a sustained increase in the general
        price level of goods and services in an economy.

        \item Inflation reduces the purchasing power of money:
        a dollar buys fewer goods and services when prices are higher.

        \item Inflation is typically measured using a price index,
        which tracks the cost of a representative basket of goods and services.

        \item The inflation rate is the percentage change in the
        price index over time:
        \[
        \pi_t = \frac{P_t-P_{t-1}}{P_{t-1}}
        \]

        \item In the U.S., common measures include the
        Consumer Price Index (CPI) and the
        Personal Consumption Expenditures (PCE) price index.
    \end{itemize}

\end{frame}

\begin{frame}[t]
    \frametitle{2. Real v. Nominal Returns}
    
    \centering
    \includegraphics[width=1\textwidth]{Econ357/figures/02_mkts_1_cpi.png}


\end{frame}

\begin{frame}[t]
    \frametitle{2. Real v. Nominal Returns}
    \framesubtitle{Definitions}

    \textbf{Definitions}

    \begin{itemize}
        \item \textbf{Nominal returns}: the percentage change in the value of an investment measured in current dollars.   Not adjusted for inflation.
        \item \textbf{Real returns}: Returns that are adjusted for inflation.   Real returns are the equivalent to the change in purchasing power.
        \begin{itemize}
            \item If prices increase, the real returns will be less than the nominal returns.   Think about an investment account that increases in value.   That's good!   You may be able to buy more stuff.  But, if prices also increase, then you will be able to buy less stuff than you otherwise would, which is less good.
        \end{itemize}
    \end{itemize}

\end{frame}

\begin{frame}[t]
    \frametitle{2. Real v. Nominal Returns}

    \centering
    \includegraphics[width=\textwidth]{Econ357/figures/02_mkts_1_1970s.png}

\end{frame}

\begin{frame}[t]
    \frametitle{2. Real v. Nominal Returns}

    \centering
    \includegraphics[width=\textwidth]{Econ357/figures/02_mkts_1_2022.png}

\end{frame}

\begin{frame}[t]
    \frametitle{2. Real v. Nominal Returns}
    \framesubtitle{Observations}

    Some observations about real returns.\\
    \begin{itemize}
      \item The appropriate price index may vary by investor
        \begin{itemize}        
        \item {\footnotesize For example, endowments may face higher price inflation if the price of health care (what they buy) increases faster than CPI.}
      \end{itemize}
      \item Often investment goals are expressed in real terms.  
      \begin{itemize}        
        \item {\footnotesize For example, "The target return for Harvard's endowed funds is 8\%, which accounts for roughly a 5\% distribution and 3\% growth to maintain purchasing power over time."}
        \item When they report returns, do they report real returns?
      \end{itemize}
      \item Straightforward to measure historical real returns.   \textbf{Expected} real returns are harder.   Specifically, it is very difficult to measure expected inflation.
    \end{itemize}
\end{frame}

\begin{frame}[t]
    \frametitle{2. Real v. Nominal Returns}

    \textbf{Be careful!}
    \vspace{1em}
    
    “Real returns” has a specific meaning – returns adjusted for inflation -- that is subtly different than the discussion of “real assets” from lecture 1.

\end{frame}

\begin{frame}[t]
    \frametitle{3. Historical Distribution of Returns}
    \section{Historical Distribution}

    The distribution of historical financial returns often looks
    approximately normal.

    \begin{itemize}
        \item The mean (or $\mu$) is the average return. This is where
        the distribution is centered.

        \item The standard deviation (or $\sigma$) measures the
        volatility of returns. It determines how widely the returns
        are dispersed around the mean.
    \end{itemize}

\end{frame}

\begin{frame}[t]
    \frametitle{3. Historical Distribution of Returns}

    \centering
    \includegraphics[width=0.9\textwidth]{Econ357/figures/02_mkts_1_sp_monthly.png}

\end{frame}

\begin{frame}[t]
    \frametitle{3. Historical Distribution of Returns}

    \centering
    \includegraphics[width=0.9\textwidth]{Econ357/figures/02_mkts_1_ust_monthly.png}

\end{frame}

\begin{frame}[t]
    \frametitle{3. Historical Distribution of Returns}

    \centering
    \includegraphics[width=0.9\textwidth]{Econ357/figures/02_mkts_1_tbills_monthly.png}

\end{frame}

\begin{frame}[t]
    \frametitle{3. Historical Distribution of Returns}
    \framesubtitle{Concerns about using the historical distribution to forecast}

    Often people use historical returns to predict expected returns.   This can be problematic.
    \begin{itemize}
        \item Financial returns often have ‘fat tails’
        \item Observations may not be independent
        \item The variance of returns may change over time
        \item Historical returns have no measure of value
    \end{itemize}
    
\end{frame}

\begin{frame}[t]
    \frametitle{3. Historical Distribution of Returns}
    \framesubtitle{Fat Tails}

    \vspace{-0.5em}

    \footnotesize
    \begin{columns}
        \begin{column}{0.5\textwidth}
                \begin{table}
                \caption{S\&P Highest and Lowest Monthly Returns Since 1950}
                \begin{tabular}{lrr}
                \toprule
                Date & Return \% & Probability \% \\
                \midrule
                 2009-04 & 12.3 & 0.054 \% \\
                1982-09 & 12.1 & 0.067 \\
                1991-02 & 11.6 & 0.109 \\
                1998-11 & 11.0 & 0.198 \\
                1975-02 & 10.8 & 0.230 \\
                \midrule
                1962-06 & -11.4 & 0.018 \% \\
                1987-10 & -11.9 & 0.011 \\
                1987-11 & -12.3 & 0.006 \\
                2020-03 & -18.9 & 0.000 \\
                2008-10 & -20.2 & 0.000 \\
                \bottomrule
                \end{tabular}
                \end{table}
        \end{column}

        \begin{column}{0.5\textwidth}
            \begin{table}
                \caption{10y UST Highest and Lowest Monthly Returns Since 1950}
                \begin{tabular}{lrr}
                \toprule
                Date & Return \% & Probability \% \\
                \midrule
                    1981-11 & 11.2 & 0.000 \% \\
                    1982-10 & 9.9 & 0.000 \\
                    2008-12 & 9.7 & 0.000 \\
                    1980-05 & 9.2 & 0.000 \\
                    1980-04 & 8.9 & 0.000 \\
                    \midrule
                    2022-09 & -5.2 & 0.150 \% \\
                    2003-07 & -5.3 & 0.120 \\
                    2022-04 & -5.4 & 0.100 \\
                    1979-10 & -5.9 & 0.040 \\
                    1980-02 & -9.6 & 0.000 \\
                \bottomrule
                \end{tabular}
            \end{table}
        \end{column}
    \end{columns}


    \blfootnote{Probabilities are based on estimated normal distribution.  
    1950-2025, n = 900 months, \n S\&P: $\mu=0.96$ and $\sigma=3.4$, 
    10y UST: $\mu=0.42$ and $\sigma=1.9$}
    
\end{frame}

\begin{frame}[t]
    \frametitle{3. Historical Distribution of Returns}
        \framesubtitle{Concerns about using the historical distribution to forecast}

    \centering
    \includegraphics[width=0.9\textwidth]{Econ357/figures/02_mkts_1_housing1.png}

\end{frame}

\begin{frame}[t]
    \frametitle{3. Historical Distribution of Returns}
    \framesubtitle{Concerns about using the historical distribution to forecast}

    \centering
    \includegraphics[width=0.9\textwidth]{Econ357/figures/02_mkts_1_housing2.png}

\end{frame}

\begin{frame}[t]
    \frametitle{3. Historical Distribution of Returns}
    \framesubtitle{What to do?}

    How financial market participants address these problems:
    \begin{enumerate}
        \item Use a distribution with fatter tails
        \item Value at Risk (VaR)
        \item Scenario Analysis / Stress Tests
    \end{enumerate}

\end{frame}

\end{document}